\documentclass[mathserif,serif,14pt,aspectratio=169]{beamer}

\usepackage{listings}
\usepackage{tikz}
\usetikzlibrary{shapes}
\usepackage{hyperref}
\usepackage{xcolor}
\usepackage{multirow}
\usepackage{eso-pic}
\usepackage[inkscapepath=ICTPresentationSvg/]{svg}
% \usepackage[cm-default]{fontspec}
% \usepackage{xeCJK}
\usepackage{ifthen}

% \setCJKmainfont{Source Han Sans CN}
% \setCJKmonofont{Source Han Sans CN}
% \setmainfont{Liberation Sans}
% \setmonofont{Courier 10 Pitch}
\usepackage{tgtermes}
\usepackage[T1]{fontenc}
\usepackage{textpos}
\usepackage{verbatim}

\def\BibTeX{{\rmfamily B\kern-.05em
\textsc{i\kern-.025em b}\kern-.08em
T\kern-.1667em\lower.7ex\hbox{E}\kern-.125emX}}

\newcommand{\link}[1]{\href{#1}{#1}}

\newcommand{\pic}[2]{
	\centerline{\includegraphics[scale=#1]{#2}}
}

% \centerline{\includegraphics[scale=#1]{#2}} % 0/0

\usepackage[backend=bibtex]{biblatex}
\addbibresource{p2p.bib}

\tikzset{grid/.style={gray,very thin,opacity=1}}	% 0/0

\definecolor{Pink}{HTML}{ff0bff}
\definecolor{Blue}{HTML}{0000ff}
\definecolor{Green}{HTML}{008a8c}
\definecolor{Brown}{HTML}{a52a2a}

\lstset{
	tabsize=4,
	columns=fixed,
	extendedchars=false,
}

\lstdefinestyle{cstyle}{
	language=c++,
	basicstyle=\ttfamily,
	morekeywords={with},
	keywordstyle=\bfseries\color{Brown},
	commentstyle=\color{Blue},
	stringstyle=\color{Pink},
	keywordstyle=[3]\color{Green},
	alsoletter={0123456789.},
	morekeywords=[4]{0,1,2,100,999},
	keywordstyle=[4]\color{Pink},
	upquote=true,
	breaklines=true,
}

\begin{document}

\AddToShipoutPictureFG*{
	\AtPageLowerLeft{
		\put(\LenToUnit{0.7\paperwidth},\LenToUnit{0.1\paperheight}){
			\begin{tikzpicture}
				\definecolor{rxx}{HTML}{ff0000}
				\definecolor{xgx}{HTML}{00ff00}
				\definecolor{xxb}{HTML}{0000ff}
				\definecolor{rgx}{HTML}{ffff00}
				\definecolor{xgb}{HTML}{00ffff}
				\definecolor{rxb}{HTML}{ff00ff}
				\definecolor{rgb}{HTML}{ffffff}
				\definecolor{xxx}{HTML}{000000}
				\definecolor{xred}{RGB}{0, 255, 255}
				\definecolor{xgreen}{RGB}{255, 0, 255}
				\filldraw [color=rxx] (0, 1) rectangle (1, 2);
				\filldraw [color=xgx] (1, 1) rectangle (2, 2);
				\filldraw [color=xxb] (2, 1) rectangle (3, 2);
				\filldraw [color=rgb] (1, 3) rectangle (2, 4);
				\filldraw [color=rgx] (0, 2) rectangle (1, 3);
				\filldraw [color=rxb] (1, 2) rectangle (2, 3);
				\filldraw [color=xgb] (2, 2) rectangle (3, 3);
				\filldraw [color=xxx] (1, 0) rectangle (2, 1);
			\end{tikzpicture}
		}
	}
}

\begin{frame}
	\centerline{\Large Implementing Synchronization Objects}
	
	\
	
	\centerline{Eric Li}
	\centerline{Feb 4, 2021}
\end{frame}

% \AtBeginSection[]{
% 	\begin{frame}
% 		\frametitle{Table of Contents}
% 		\tableofcontents[currentsection]
% 	\end{frame}
% }

\section{Content}

\begin{frame}
	\frametitle{Overview}
	\begin{itemize}\setlength\itemsep{0.5em}
	\item<1-> Goal: implement locks and condition variables (CV)
	\item<2-> We need to atomically modify states
		\begin{itemize}\setlength\itemsep{0.3em}
		\item<2-> Lock: FREE / BUSY; queue of waiting threads
		\item<2-> CV: queue of waiting threads
		\end{itemize}
	\item<3-> Method: use hardware primitives
		\begin{itemize}\setlength\itemsep{0.3em}
		\item<3-> Disable interrupt (uniprocessor)
		\item<3-> Atomic read-modify-write instructions (multiprocessor)
		\end{itemize}
	\item<4-> Kernel mode vs. user mode
	\end{itemize}
\end{frame}

\begin{frame} [fragile]
	\frametitle{Uniprocessor Locks: Disabling Interrupts}
	\begin{itemize}\setlength\itemsep{0.5em}
	\item<1-> No other processors can change memory
	\item<1-> If no context switch, operations appear to be atomic
	\item<2-> Trivial implementation
		{
		\small
		\begin{lstlisting}[style=cstyle, gobble=8, frame=single, 
							numbers=left, numberstyle=\ttfamily]
		Lock::acquire() { disableInterrupts(); }
		Lock::release() { enableInterrupts(); }
		\end{lstlisting}
		}
	\item<3-> Problem
		\begin{itemize}\setlength\itemsep{0.3em}
		\item<3-> Disable interrupt for a long time (starvation, not real-time)
		\item<3-> Cannot allow user-level code to disable interrupts
		\end{itemize}
	\end{itemize}
\end{frame}

\begin{frame}
	\frametitle{Implementing Uniprocessor Queueing Locks}
	\begin{itemize}\setlength\itemsep{0.5em}
	\item<1-> Briefly disable interrupt to protect lock structure
	\item<1-> When lock is locked, context switch to another ready thread
	\end{itemize}
\end{frame}

\begin{frame} [fragile]
	\frametitle{Implementing Uniprocessor Queueing Locks}
	{
	\fontsize{8}{8}\selectfont
	\begin{columns}
		\begin{column}{0.50\textwidth}
		\begin{lstlisting}[style=cstyle, gobble=8, frame=single, 
							numbers=left, numberstyle=\ttfamily]
		Lock::acquire() {
			TCB *chosenTCB;

			disableInterrupts();
			if (value == BUSY) {
				waiting.add(runningThread);
				runningThread->state = WAITING;
				chosenTCB = readyList.remove();
				thread_switch(runningThread,
				              chosenTCB);
				runningThread->state = RUNNING;
			} else {
				value = BUSY;
			}
			enableInterrupts();
		}
		\end{lstlisting}
		\end{column}
		\begin{column}{0.45\textwidth}
		\begin{lstlisting}[style=cstyle, gobble=8, frame=single, 
							numbers=left, numberstyle=\ttfamily]
		Lock::release() {
		// next thread to hold lock
			TCB *next;

			disableInterrupts();
			if (waiting.notEmpty()) {
			// move one TCB from waiting
			// to ready
				next = waiting.remove();
				next->state = READY;
				readyList.add(next);
			} else {
				value = FREE;
			}
			enableInterrupts();
		}
		\end{lstlisting}
		\end{column}
	\end{columns}
	}
	\begin{itemize}\setlength\itemsep{0.5em}
	\item<2-> Do not set value = FREE in release() to prevent starvation
		\begin{textblock}{2}[0,0](8,-4.5)
			{\color{red} \textbf{-->}}	% TODO
		\end{textblock}
	\item<3-> During call to \lstinline{thread_switch}, interrupts are turned off
		\begin{textblock}{2}[0,0](0.7,-7.2)
			{\color{red} \textbf{-->}}	% TODO
		\end{textblock}
	\end{itemize}
\end{frame}

\begin{frame} [fragile]
	\frametitle{Implementing Multiprocessor Spinlocks}
	\begin{itemize}\setlength\itemsep{0.5em}
	\item<1-> Cannot turn off interrupt
	\item<1-> Use atomic read-modify-write instructions
		\begin{itemize}\setlength\itemsep{0.3em}
		\item<1-> Implementation: related to cache
		\item<1-> Computer Architecture courses (ECS 154B / probably ECS 201)
		\item<1->
			\href{https://people.freebsd.org/~lstewart/articles/cpumemory.pdf
				}{"What Every Programmer Should Know About Memory"}
				Figure 3.18
		\end{itemize}
	\item<2-> Atomic test-and-set instruction
		{
		\small
		\begin{lstlisting}[style=cstyle, gobble=8, frame=single, 
							numbers=left, numberstyle=\ttfamily]
		int test_and_set(int* lockPtr, int newValue) {
			int oldValue;
			oldValue = *lockPtr;
			*lockPtr = newValue;
			return oldValue;
		}
		\end{lstlisting}
		}
	\end{itemize}
\end{frame}

\begin{frame} [fragile]
	\frametitle{Implementing Multiprocessor Spinlocks}
	{
	\fontsize{8}{8}\selectfont
	\begin{lstlisting}[style=cstyle, gobble=4, frame=single, 
						numbers=left, numberstyle=\ttfamily]
	class SpinLock {
	  private:
		int value = 0; // 0 = FREE; 1 = BUSY

	  public:
		void acquire() {
			while (test_and_set(&value)) // while BUSY
			    ; // spin
		}

		void release() {
			value = 0;
			memory_barrier();
		}
	}
	\end{lstlisting}
	}
	\begin{itemize}\setlength\itemsep{0.5em}
	\item<2-> Busy wait (assume locks are only held shortly)
		\begin{textblock}{2}[0,0](9.2,-6.4)
			{\color{red} \textbf{<--}}	% TODO
		\end{textblock}
	\end{itemize}
\end{frame}

\begin{frame} [fragile]
	\frametitle{Implementing Multiprocessor Queueing Locks}
	\begin{itemize}\setlength\itemsep{0.5em}
	\item<1-> Critical section length can be long
	\item<1-> Minimize busy waiting
	\item<2-> Class definitions
		{
		\fontsize{8}{8}\selectfont
		\begin{lstlisting}[style=cstyle, gobble=8, frame=single, 
							numbers=left, numberstyle=\ttfamily]
		class Lock {
		  private:
			int value = FREE;
			SpinLock spinLock;
			Queue waiting;
		  public:
			void acquire();
			void release();
		}
		class Scheduler {
		  private:
			Queue readyList;
			SpinLock schedulerSpinLock;
		  public:
			void suspend(SpinLock *lock);
			void makeReady(Thread *thread);
		}
		\end{lstlisting}
		}
	\end{itemize}
\end{frame}

\begin{frame} [fragile]
	\frametitle{Implementing Uniprocessor Queueing Locks}
	{
	\fontsize{7}{7}\selectfont
	\begin{columns}
		\begin{column}{0.47\textwidth}
		\lstinputlisting[style=cstyle, frame=single, 
						numbers=left, numberstyle=\ttfamily]{574a.cpp}
		\end{column}
		\begin{column}{0.49\textwidth}
		\lstinputlisting[style=cstyle, frame=single, 
						numbers=left, numberstyle=\ttfamily]{574b.cpp}
		\end{column}
	\end{columns}
	}
	{
	\fontsize{12}{12}\selectfont
	\begin{itemize}
	\item<2-> Suspending a thread
		\begin{enumerate}
		\fontsize{10}{10}\selectfont
		\item<3-> Call scheduler.suspend without releasing Lock's spinLock
			\begin{textblock}{2}[0,0](0.7,-10)
				{\color{red} \textbf{-->}}	% TODO
			\end{textblock}
		\item<4-> disableInterrupts() to prevent thread being preempted
			\begin{textblock}{2}[0,0](7.5,-11.6)
				{\color{red} \textbf{-->}}	% TODO
			\end{textblock}
		\item<5-> Acquire Scheduler's spinLock to protect readyList
			\begin{textblock}{2}[0,0](7.5,-11.7)
				{\color{red} \textbf{-->}}	% TODO
			\end{textblock}
		\end{enumerate}
	\end{itemize}
	}
\end{frame}

\begin{frame} [fragile]
	\frametitle{Implementing Uniprocessor Queueing Locks}
	{
	\fontsize{7}{7}\selectfont
	\begin{columns}
		\begin{column}{0.47\textwidth}
		\lstinputlisting[style=cstyle, frame=single, 
						numbers=left, numberstyle=\ttfamily]{574c.cpp}
		\end{column}
		\begin{column}{0.49\textwidth}
		\lstinputlisting[style=cstyle, frame=single, 
						numbers=left, numberstyle=\ttfamily]{574d.cpp}
		\end{column}
	\end{columns}
	}
	{
	\fontsize{12}{12}\selectfont
	\begin{itemize}
	\item<1-> What if we release Lock's spinLock before calling scheduler.suspend?
		\begin{enumerate}
		\fontsize{10}{10}\selectfont
		\item<2-> After release spinLock
			\begin{textblock}{2}[0,0](0.7,-9.9)
				{\color{red} \textbf{-->}}	% TODO
			\end{textblock}
		\item<3-> Another thread release lock, make this thread READY
			\begin{textblock}{2}[0,0](7.5,-4.5)
				{\color{red} \textbf{-->}}	% TODO
			\end{textblock}
		\item<4-> Current thread calls suspend(), state = WAITING forever
			\begin{textblock}{2}[0,0](7.5,-10.8)
				{\color{red} \textbf{-->}}	% TODO
			\end{textblock}
		\end{enumerate}
	\end{itemize}
	}
\end{frame}

\begin{frame} [fragile]
	\frametitle{Linux 2.6 Kernel Mutex Lock}
	\begin{itemize}\setlength\itemsep{0.5em}
	\item<1-> Optimized for the common case
	\item<1-> Assumption: most locks are FREE most of the time
	\item<2-> Acquire: fast path when lock is not already acquired
	\item<3-> Release: fast path when no waiters on the lock
	\end{itemize}
\end{frame}

\begin{frame} [fragile]
	\frametitle{Implementing Condition Variables}
	\begin{itemize}\setlength\itemsep{0.5em}
	\item<1-> Similar to implementing locks
	\item<2-> Still, pass spinLock to scheduler.suspend
		\begin{textblock}{2}[0,0](5.2,6.9)
			{\color{red} \textbf{<--}}	% TODO
		\end{textblock}
	\end{itemize}
	{
	\fontsize{8}{8}\selectfont
	\begin{columns}
		\begin{column}{0.47\textwidth}
		\begin{lstlisting}[style=cstyle, gobble=8, frame=single, 
							numbers=left, numberstyle=\ttfamily]
		class CV {
		  private:
			Queue waiting;
		  public:
			void wait(Lock *lock);
			void signal();
			void broadcast();
		}

		void CV::wait(Lock *lock) {
			assert(lock.isHeld());
			waiting.add(myTCB);
			scheduler.suspend(&lock);
			lock->acquire();
		}
		\end{lstlisting}
		\end{column}
		\begin{column}{0.47\textwidth}
		\begin{lstlisting}[style=cstyle, gobble=8, frame=single, 
							numbers=left, numberstyle=\ttfamily]
		void CV::signal() {
			if (waiting.notEmpty()) {
				thread = waiting.remove();
				scheduler.makeReady(thread);
			}
		}

		void CV::broadcast() {
			while (waiting.notEmpty()) {
				thread = waiting.remove();
				scheduler.makeReady(thread);
			}
		}
		\end{lstlisting}
		\end{column}
	\end{columns}
	}
\end{frame}

\begin{frame} [fragile]
	\frametitle{Implementing Application-level Synchronization}
	\begin{itemize}\setlength\itemsep{0.5em}
	\item<1-> Kernel-Managed Threads
		\begin{itemize}\setlength\itemsep{0.5em}
		\item<1-> Simple case: place Lock and CV in kernel space, app use syscalls
		\item<2-> Sophisticated case: fast path in user space, slow path in kernel
		\end{itemize}
	\item<3-> User-Managed Threads
		\begin{itemize}\setlength\itemsep{0.5em}
		\item Implement most things at user level
		\item Disabling interrupts $\to$ temporarily disable upcalls
			(usually supported by modern OS)
		\end{itemize}
	\end{itemize}
\end{frame}

\section{EOF}

\AddToShipoutPictureFG*{
	\AtPageLowerLeft{
		\put(\LenToUnit{0.05\paperwidth},\LenToUnit{0.26\paperheight}){
			Ref: Anderson \& Dahlin, Operating Systems - Principles and Practice
		}
		\put(\LenToUnit{0.05\paperwidth},\LenToUnit{0.18\paperheight}){
			Ref: \url{https://en.wikipedia.org/wiki/Test-and-set}
		}
		\put(\LenToUnit{0.05\paperwidth},\LenToUnit{0.1\paperheight}){
			Thanks: \LaTeX, Beamer, OBS
		}
	}
}

\begin{frame}
	\centerline{\fontsize{32}{32}\selectfont Thank you}
\end{frame}

% Words: anonymity, reciprocity, proprietary

\newpage
\end{document}

